\documentclass[a4paper, serif, 10pt]{moderncv}

%% moderncv
% color and style
\moderncvcolor{black}
\moderncvstyle{banking}

%% page layout/margins
\usepackage[top=2.5cm, bottom=2.5cm, left=1.5cm, right=1.5cm]{geometry}

\nopagenumbers{}

%% Babel config for French language
% ref:
% - https://latex3.github.io/babel/guides/locale-french.html
\usepackage[french]{babel}

%% fontspec
\usepackage{fontspec}

\IfFontExistsTF{Linux Libertine}{
  \setmainfont[Ligatures={TeX, Common}, Numbers=OldStyle]{Linux Libertine}
}{}
\IfFontExistsTF{Linux Libertine O}{
  \setmainfont[Ligatures={TeX, Common}, Numbers=OldStyle]{Linux Libertine O}
}{}

%% CTeX
% CJK support, used to show CJK characters in the resume
%
% - heading=false: disable ctex heading
% - punct=kaiming: use kaiming punctuations styles for CJK
% - scheme=plain: use plain scheme, do not override \normalsize font size
% - space=auto: space settings for CJK characters
%
% CJK fonts are resolved per locale language, e.g. Noto Serif/Sans CJK JP
% for Japanese resumes, with fallback to the next region if not installed
%
% ref:
% - http://ctan.mirrorcatalogs.com/language/chinese/ctex/ctex.pdf
\usepackage[UTF8, heading=false, punct=kaiming, scheme=plain, space=auto]{ctex}

\IfFontExistsTF{Noto Serif CJK SC}{
  \setCJKmainfont{Noto Serif CJK SC}
}{}
\IfFontExistsTF{Noto Sans CJK SC}{
  \setCJKsansfont{Noto Sans CJK SC}
}{}

%% line spacing
\usepackage{setspace}
\setstretch{1.125}

%% URL styling - use normal text instead of monospace
\urlstyle{same}

%% Patch \httplink and \httpslink to handle full URLs
% The original moderncv macros always prepend http:// or https:// to URLs.
% This causes issues when the URL already has a protocol (e.g., https://example.com)
% resulting in malformed URLs like https://https://example.com.
% This patch checks if the URL already contains "://" and uses it directly if so.
% Note: We use \str_set:Nx (x-expansion) to expand macros like \@homepage first.
\ExplSyntaxOn
\str_new:N \g_yamlresume_url_str

\RenewDocumentCommand{\httpslink}{O{}m}{
  \str_set:Nx \g_yamlresume_url_str {#2}
  \str_if_in:NnTF \g_yamlresume_url_str {://}
    { \url{#2} }
    { \url{https://#2} }
}

\RenewDocumentCommand{\httplink}{O{}m}{
  \str_set:Nx \g_yamlresume_url_str {#2}
  \str_if_in:NnTF \g_yamlresume_url_str {://}
    { \url{#2} }
    { \url{http://#2} }
}
\ExplSyntaxOff

\name{Camille Lefebvre}{}
\title{Analyste d'Affaires Senior | Expert en Transformation Digitale}
\phone[mobile]{+33 6 12 34 56 78}
\email{camille.lefebvre@email.fr}
\homepage{https://www.camille-lefebvre.fr}

\address{25 Rue de la Paix, Paris, Île-de-France, France, 75002}{}{}

\extrainfo{{\small \faLinkedin}\ \href{https://www.linkedin.com/in/camille-lefebvre}{@camille-lefebvre} {} {} {} • {} {} {} 
{\small \faGithub}\ \href{https://github.com/camillelefebvre}{@camillelefebvre}}

\begin{document}

\maketitle

\section{Informations de base}

\cvline{}{\begin{itemize}
\item Analyste d'affaires avec 8 ans d'expérience dans la conduite de projets de transformation digitale.
\item Expertise en modélisation des processus métier (BPMN, UML) et en analyse fonctionnelle.
\item Maîtrise des outils d'analyse de données (SQL, Power BI, Tableau) et des méthodologies agiles.
\item Forte capacité à collaborer avec les équipes techniques et métier pour aligner les solutions sur les besoins business.
\end{itemize}}
\section{Formation}

\cventry{sept. 2014–juil. 2016}
        {Master, Analyse d'Affaires et Systèmes d'Information, Score : Mention Très Bien}
        {Université Paris-Dauphine}
        {\url{https://www.dauphine.psl.eu}}
        {}
        {\begin{itemize}
\item Formation approfondie en analyse fonctionnelle et modélisation des processus.
\item Projets pratiques en collaboration avec des entreprises partenaires.
\item Mémoire sur l'impact de la transformation digitale sur les PME.
\end{itemize}
\textbf{Cours} : Modélisation des processus métier, Gestion de projet agile, Systèmes d'information décisionnels, Analyse de données, Architecture des SI, Management de l'innovation, Droit des affaires, Communication professionnelle}
\section{Expérience professionnelle}

\cventry{janv. 2021–Present}
        {Analyste d'Affaires Senior}
        {BNP Paribas}
        {\url{https://www.bnpparibas.fr}}
        {}
        {\begin{itemize}
\item Pilotage de l'analyse fonctionnelle pour le projet de refonte du portail client, en collaboration avec les équipes métier et IT.
\item Rédaction des spécifications fonctionnelles détaillées (SFG) et des cahiers de recette.
\item Animation d'ateliers de cadrage et de priorisation avec les parties prenantes.
\item Modélisation des processus métier en BPMN et création de maquettes fonctionnelles.
\item Suivi des indicateurs de performance et reporting mensuel auprès du comité de direction.
\end{itemize}
\textbf{Mots-clés} : BPMN, Analyse fonctionnelle, Agile, JIRA, Confluence}

\cventry{sept. 2016–déc. 2020}
        {Analyste d'Affaires}
        {Capgemini}
        {\url{https://www.capgemini.com/fr-fr/}}
        {}
        {\begin{itemize}
\item Participation à des missions de conseil en transformation digitale pour des clients du secteur bancaire et de l'assurance.
\item Analyse des besoins métier et rédaction des user stories.
\item Coordination entre les équipes de développement et les utilisateurs finaux.
\item Réalisation de tests fonctionnels et validation des livrables.
\item Contribution à l'amélioration continue des processus internes.
\end{itemize}
\textbf{Mots-clés} : User stories, UML, Tests fonctionnels, Scrum, SQL}
\section{Langues}

\cvline{Français}{Niveau natif ou bilingue}
\cvline{Anglais}{Niveau professionnel complet \hfill \textbf{Mots-clés} : TOEIC 950}
\cvline{Espagnol}{Niveau de travail limité}
\section{Compétences}

\cvline{Analyse et Modélisation}{Expert \hfill \textbf{Mots-clés} : BPMN, UML, Analyse fonctionnelle, Cartographie des processus}
\cvline{Gestion de Projet}{Avancé \hfill \textbf{Mots-clés} : Agile, Scrum, JIRA, Confluence, MS Project}
\cvline{Analyse de Données}{Avancé \hfill \textbf{Mots-clés} : SQL, Power BI, Tableau, Excel avancé}
\cvline{Outils Bureautiques}{Expert \hfill \textbf{Mots-clés} : Excel, PowerPoint, Word, Visio}
\section{Récompenses}

\cventry{juin 2019}
        {Capgemini}
        {Prix de l'Innovation}
        {}
        {}
        {Récompense pour la conception d'un outil interne d'analyse des besoins clients.}
\section{Certificats}

\cventry{mars 2022}
        {IIBA}
        {Certified Business Analysis Professional (CBAP)}
        {\url{https://www.iiba.org/certification/cbap/}}
        {}
        {}

\cventry{nov. 2020}
        {Scrum.org}
        {Professional Scrum Master I}
        {\url{https://www.scrum.org/professional-scrum-master-i-certification}}
        {}
        {}
\section{Publications}

\cventry{avr. 2021}
        {L'analyse d'affaires à l'ère du digital : bonnes pratiques et outils}
        {Journal du Management}
        {\url{https://www.journaldumanagement.fr/analyse-affaires-digital}}
        {}
        {Article présentant les méthodes et outils clés pour réussir la transformation digitale des organisations.}
\section{Projets}

\cventry{janv. 2021–déc. 2022}
        {Analyse des besoins, rédaction des spécifications fonctionnelles et coordination avec les équipes de développement pour la refonte du portail client.
}
        {Refonte du portail client}
        {\url{https://www.bnpparibas.fr/portail-client}}
        {}
        {Projet de modernisation du portail client de BNP Paribas.
\textbf{Mots-clés} : BPMN, Agile, UX}

\cventry{mars 2020–août 2020}
        {Conception et mise en place d'un tableau de bord interactif sous Power BI pour visualiser les KPIs métier et faciliter la prise de décision.
}
        {Outil de pilotage des indicateurs}
        {}
        {}
        {Développement d'un tableau de bord Power BI pour le suivi des indicateurs de performance.
\textbf{Mots-clés} : Power BI, SQL, KPIs}
\section{Centres d'intérêt}

\cvline{Lecture}{Romans, Essais, Biographies}
\cvline{Randonnée}{Montagne, Nature}
\cvline{Photographie}{Voyage, Portrait}
\section{Bénévolat}

\cventry{sept. 2020–Present}
        {Mentore}
        {Association Femmes du Numérique}
        {\url{https://www.femmesdunumerique.fr}}
        {}
        {Accompagnement de jeunes femmes dans leur orientation vers les métiers du numérique, animation d'ateliers et partage d'expérience.}
    
\end{document}